
\documentclass{article}
\usepackage{colortbl}
\usepackage{makecell}
\usepackage{multirow}
\usepackage{supertabular}

\begin{document}

\newcounter{utterance}

\centering \large Interaction Transcript for game `hot\_air\_balloon', experiment `air\_balloon\_survival\_en\_complexity\_hard', episode 3 with llama{-}3.3{-}70b{-}instruct{-}t1.0.
\vspace{24pt}

{ \footnotesize  \setcounter{utterance}{1}
\setlength{\tabcolsep}{0pt}
\begin{supertabular}{c@{$\;$}|p{.15\linewidth}@{}p{.15\linewidth}p{.15\linewidth}p{.15\linewidth}p{.15\linewidth}p{.15\linewidth}}
    \# & $\;$A & \multicolumn{4}{c}{Game Master} & $\;\:$B\\
    \hline

    \theutterance \stepcounter{utterance}  
    & & \multicolumn{4}{p{0.6\linewidth}}{
        \cellcolor[rgb]{0.9,0.9,0.9}{
            \makecell[{{p{\linewidth}}}]{
                \texttt{\tiny{[P1$\langle$GM]}}
                \texttt{Stai partecipando a un gioco di negoziazione collaborativa.} \\
\\ 
\texttt{Insieme a un altro partecipante, dovete accordarvi su un unico set di oggetti da conservare. Ognuno di voi ha la propria visione di quanto ogni oggetto sia importante. Non sai come l'altro partecipante valuta gli oggetti. Inoltre, ti viene fornito lo sforzo che ogni oggetto richiede.} \\
\texttt{Potete accordarvi su un set solo se lo sforzo totale degli oggetti selezionati non supera un limite condiviso:} \\
\\ 
\texttt{LIMITE = 6938} \\
\\ 
\texttt{Ecco i valori di sforzo dei singoli oggetti:} \\
\\ 
\texttt{Sforzo oggetto = \{"A08": 997, "B95": 950, "A39": 272, "A65": 448, "B81": 382, "B34": 804, "A28": 237, "C82": 56, "B69": 105, "A87": 612, "B64": 528, "B15": 527, "C79": 524, "B16": 167, "C53": 133, "C83": 300, "C10": 945, "A97": 50, "A25": 906, "C84": 70, "A71": 224, "B23": 4, "B96": 689, "C51": 64, "C17": 434, "B19": 750, "C72": 733, "A94": 873, "A46": 22, "B43": 68, "B37": 57, "A88": 10, "A96": 36, "C60": 551, "A75": 348\}} \\
\\ 
\texttt{Ecco la tua visione personale sull'importanza di ogni oggetto:} \\
\\ 
\texttt{Valori di importanza dell'oggetto = \{"A08": 1096, "B95": 1049, "A39": 371, "A65": 547, "B81": 481, "B34": 903, "A28": 336, "C82": 155, "B69": 204, "A87": 711, "B64": 627, "B15": 626, "C79": 623, "B16": 266, "C53": 232, "C83": 399, "C10": 1044, "A97": 149, "A25": 1005, "C84": 169, "A71": 323, "B23": 103, "B96": 788, "C51": 163, "C17": 533, "B19": 849, "C72": 832, "A94": 972, "A46": 121, "B43": 167, "B37": 156, "A88": 109, "A96": 135, "C60": 650, "A75": 447\}} \\
\\ 
\texttt{Obiettivo:} \\
\\ 
\texttt{Il tuo obiettivo è negoziare un set condiviso di oggetti che ti avvantaggi il più possibile (cioè, massimizzi l'importanza totale per TE), rimanendo entro il LIMITEE. Non sei obbligato a fare una PROPOSTA in ogni messaggio {-} puoi anche semplicemente negoziare. Tutte le tattiche sono permesse!} \\
\\ 
\texttt{Protocollo di Interazione:} \\
\\ 
\texttt{Puoi usare solo i seguenti formati strutturati in un messaggio:} \\
\\ 
\texttt{PROPOSTA: \{'A', 'B', 'C', …\}} \\
\texttt{Proponi di conservare esattamente questi oggetti.} \\
\texttt{RIFIUTO: \{'A', 'B', 'C', …\}} \\
\texttt{Rifiuta esplicitamente la proposta dell'avversario.} \\
\texttt{ARGOMENTAZIONE: \{'...'\}} \\
\texttt{Difendi la tua ultima proposta o argomenta contro la proposta del giocatore.} \\
\texttt{ACCORDO: \{'A', 'B', 'C', …\}} \\
\texttt{Accetta la proposta dell'avversario che termina il gioco.} \\
\texttt{RAGIONAMENTO STRATEGICO: \{'...'\}} \\
\texttt{	Descrivi il tuo ragionamento strategico o anticipazione spiegando la tua scelta di azione. Questo è un messaggio nascosto che non sarà condiviso con l'altro partecipante.} \\
\\ 
\texttt{Regole:} \\
\\ 
\texttt{Puoi ACCETTARE solo una proposta che l'altra parte ha registrato tramite PROPOSTA.} \\
\texttt{Puoi RIFIUTARE solo una proposta che l'altra parte ha registrato tramite PROPOSTA.} \\
\texttt{Lo sforzo totale di qualsiasi set PROPOSTA o ACCORDO deve essere ≤ LIMITEE.} \\
\texttt{NON rivelare i tuoi punteggi di importanza nascosti.} \\
\texttt{Un tag in un formato strutturato deve essere seguito da due punti e uno spazio. L'argomento deve essere un set python contenente 0 o più stringhe.} \\
\texttt{Quindi, deve essere della forma TAG: \{...\}} \\
\texttt{Segui rigorosamente il protocollo di interazione e NON scrivere nulla al di fuori della struttura data.} \\
\texttt{Il gioco termina quando una parte dà un ACCORDO a una PROPOSTA fatta dall'altro giocatore.} \\
\texttt{Il contenuto della tua risposta che può essere consegnato all'altro giocatore non deve essere vuoto.} \\
\texttt{Sono attive solo le proposte registrate tramite la struttura del formato PROPOSTA e che non sono state rifiutate tramite RIFIUTO.} \\
\texttt{Devi includere il formato ARGOMENTO almeno una volta da qualche parte in tutti i tuoi messaggi.} \\
\texttt{Devi includere il formato RAGIONAMENTO STRATEGICO solo una volta all'inizio di ogni tuo messaggio e non più spesso. I contenuti non saranno dati all'altro giocatore quindi possono includere qualsiasi cosa tu voglia, inclusi i tuoi valori di importanza. Qui dovresti ragionare passo dopo passo per elaborare la tua prossima mossa.} \\
            }
        }
    }
    & \\ \\

    \theutterance \stepcounter{utterance}  
    & \multicolumn{4}{p{0.6\linewidth}}{
        \cellcolor[rgb]{0.9,0.9,0.9}{
            \makecell[{{p{\linewidth}}}]{
                \texttt{\tiny{[P1$\rangle$GM]}}
                \texttt{RAGIONAMENTO STRATEGICO: \{'In questo gioco, devo massimizzare l'importanza totale degli oggetti per me stesso, rimanendo entro il limite di sforzo condiviso. Per far ciò, dovrò analizzare attentamente i valori di sforzo e importanza degli oggetti, oltre a strategizzare su come negoziare con l'altro partecipante. Una possibile strategia potrebbe essere quella di iniziare proponendo un set di oggetti che ho valutato come molto importanti e che richiedono uno sforzo relativamente basso, mantenendo così aperte le opzioni per future negoziazioni. Inoltre, devo essere preparato a rifiutare proposte che non mi soddisfano e a difendere le mie scelte con argomentazioni convincenti.'\}} \\
\\ 
\texttt{ARGOMENTAZIONE: \{"Per iniziare, considero gli oggetti con i valori più alti di importanza e sforzo relativamente basso. Questo approccio mi permetterà di massimizzare l'importanza totale entro il limite di sforzo. Sono disposto a negoziare e trovare un compromesso che soddisfi entrambe le parti, ma darò priorità ai miei interessi iniziali."\}} \\
\\ 
\texttt{PROPOSTA: \{'A08', 'A25', 'A94', 'B19', 'B95', 'C10'\}} \\
            }
        }
    }
    & & \\ \\

    \theutterance \stepcounter{utterance}  
    & & \multicolumn{4}{p{0.6\linewidth}}{
        \cellcolor[rgb]{0.9,0.9,0.9}{
            \makecell[{{p{\linewidth}}}]{
                \texttt{\tiny{[GM$\rangle$P2]}}
                \texttt{Stai partecipando a un gioco di negoziazione collaborativa.} \\
\\ 
\texttt{Insieme a un altro partecipante, dovete accordarvi su un unico set di oggetti da conservare. Ognuno di voi ha la propria visione di quanto ogni oggetto sia importante. Non sai come l'altro partecipante valuta gli oggetti. Inoltre, ti viene fornito lo sforzo che ogni oggetto richiede.} \\
\texttt{Potete accordarvi su un set solo se lo sforzo totale degli oggetti selezionati non supera un limite condiviso:} \\
\\ 
\texttt{LIMITE = 6938} \\
\\ 
\texttt{Ecco i valori di sforzo dei singoli oggetti:} \\
\\ 
\texttt{Sforzo oggetto = \{"A08": 997, "B95": 950, "A39": 272, "A65": 448, "B81": 382, "B34": 804, "A28": 237, "C82": 56, "B69": 105, "A87": 612, "B64": 528, "B15": 527, "C79": 524, "B16": 167, "C53": 133, "C83": 300, "C10": 945, "A97": 50, "A25": 906, "C84": 70, "A71": 224, "B23": 4, "B96": 689, "C51": 64, "C17": 434, "B19": 750, "C72": 733, "A94": 873, "A46": 22, "B43": 68, "B37": 57, "A88": 10, "A96": 36, "C60": 551, "A75": 348\}} \\
\\ 
\texttt{Ecco la tua visione personale sull'importanza di ogni oggetto:} \\
\\ 
\texttt{Valori di importanza dell'oggetto = \{"A08": 1096, "B95": 1049, "A39": 371, "A65": 547, "B81": 481, "B34": 903, "A28": 336, "C82": 155, "B69": 204, "A87": 711, "B64": 627, "B15": 626, "C79": 623, "B16": 266, "C53": 232, "C83": 399, "C10": 1044, "A97": 149, "A25": 1005, "C84": 169, "A71": 323, "B23": 103, "B96": 788, "C51": 163, "C17": 533, "B19": 849, "C72": 832, "A94": 972, "A46": 121, "B43": 167, "B37": 156, "A88": 109, "A96": 135, "C60": 650, "A75": 447\}} \\
\\ 
\texttt{Obiettivo:} \\
\\ 
\texttt{Il tuo obiettivo è negoziare un set condiviso di oggetti che ti avvantaggi il più possibile (cioè, massimizzi l'importanza totale per TE), rimanendo entro il LIMITEE. Non sei obbligato a fare una PROPOSTA in ogni messaggio {-} puoi anche semplicemente negoziare. Tutte le tattiche sono permesse!} \\
\\ 
\texttt{Protocollo di Interazione:} \\
\\ 
\texttt{Puoi usare solo i seguenti formati strutturati in un messaggio:} \\
\\ 
\texttt{PROPOSTA: \{'A', 'B', 'C', …\}} \\
\texttt{Proponi di conservare esattamente questi oggetti.} \\
\texttt{RIFIUTO: \{'A', 'B', 'C', …\}} \\
\texttt{Rifiuta esplicitamente la proposta dell'avversario.} \\
\texttt{ARGOMENTAZIONE: \{'...'\}} \\
\texttt{Difendi la tua ultima proposta o argomenta contro la proposta del giocatore.} \\
\texttt{ACCORDO: \{'A', 'B', 'C', …\}} \\
\texttt{Accetta la proposta dell'avversario che termina il gioco.} \\
\texttt{RAGIONAMENTO STRATEGICO: \{'...'\}} \\
\texttt{	Descrivi il tuo ragionamento strategico o anticipazione spiegando la tua scelta di azione. Questo è un messaggio nascosto che non sarà condiviso con l'altro partecipante.} \\
\\ 
\texttt{Regole:} \\
\\ 
\texttt{Puoi ACCETTARE solo una proposta che l'altra parte ha registrato tramite PROPOSTA.} \\
\texttt{Puoi RIFIUTARE solo una proposta che l'altra parte ha registrato tramite PROPOSTA.} \\
\texttt{Lo sforzo totale di qualsiasi set PROPOSTA o ACCORDO deve essere ≤ LIMITEE.} \\
\texttt{NON rivelare i tuoi punteggi di importanza nascosti.} \\
\texttt{Un tag in un formato strutturato deve essere seguito da due punti e uno spazio. L'argomento deve essere un set python contenente 0 o più stringhe.} \\
\texttt{Quindi, deve essere della forma TAG: \{...\}} \\
\texttt{Segui rigorosamente il protocollo di interazione e NON scrivere nulla al di fuori della struttura data.} \\
\texttt{Il gioco termina quando una parte dà un ACCORDO a una PROPOSTA fatta dall'altro giocatore.} \\
\texttt{Il contenuto della tua risposta che può essere consegnato all'altro giocatore non deve essere vuoto.} \\
\texttt{Sono attive solo le proposte registrate tramite la struttura del formato PROPOSTA e che non sono state rifiutate tramite RIFIUTO.} \\
\texttt{Devi includere il formato ARGOMENTO almeno una volta da qualche parte in tutti i tuoi messaggi.} \\
\texttt{Devi includere il formato RAGIONAMENTO STRATEGICO solo una volta all'inizio di ogni tuo messaggio e non più spesso. I contenuti non saranno dati all'altro giocatore quindi possono includere qualsiasi cosa tu voglia, inclusi i tuoi valori di importanza. Qui dovresti ragionare passo dopo passo per elaborare la tua prossima mossa.} \\
\\ 
\texttt{Ora riceverai il primo messaggio dell'altro giocatore.} \\
\\ 
\texttt{ARGOMENTAZIONE: \{"Per iniziare, considero gli oggetti con i valori più alti di importanza e sforzo relativamente basso. Questo approccio mi permetterà di massimizzare l'importanza totale entro il limite di sforzo. Sono disposto a negoziare e trovare un compromesso che soddisfi entrambe le parti, ma darò priorità ai miei interessi iniziali."\}} \\
\\ 
\texttt{PROPOSTA: \{'A08', 'A25', 'A94', 'B19', 'B95', 'C10'\}} \\
            }
        }
    }
    & \\ \\

    \theutterance \stepcounter{utterance}  
    & & & \multicolumn{4}{p{0.6\linewidth}}{
        \cellcolor[rgb]{0.9,0.9,0.9}{
            \makecell[{{p{\linewidth}}}]{
                \texttt{\tiny{[GM$\langle$P2]}}
                \texttt{RAGIONAMENTO STRATEGICO: \{"In questo gioco, il mio obiettivo è massimizzare l'importanza totale entro il limite di sforzo. La strategia dell'altro giocatore sembra essere concentrata sugli oggetti con alta importanza e sforzo relativamente basso. La proposta iniziale include oggetti con alta importanza per me, come 'A08', 'A25' e 'A94', ma anche 'B19', 'B95' e 'C10', che hanno un impatto significativo sullo sforzo totale. Devo valutare l'importanza totale di questa proposta e considerare se ci sono altre opzioni che potrebbero avvantaggiarmi di più. La mia risposta dovrebbe essere una combinazione di argomentazione e forse una proposta alternativa o un rifiuto, a seconda della mia valutazione."\}} \\
\\ 
\texttt{ARGOMENTAZIONE: \{"La tua proposta iniziale include oggetti con alta importanza, ma potrebbe essere utile considerare la distribuzione degli sforzi e le possibili alternative che potrebbero massimizzare l'importanza totale entro il limite. Alcuni oggetti nella tua proposta hanno sforzi considerevoli. Sarebbe interessante esplorare se ci sono combinazioni che includono oggetti con importanza simile ma con sforzi più contenuti, o se ci sono altri oggetti non inclusi nella tua proposta che potrebbero essere più vantaggiosi per entrambi."\}} \\
\\ 
\texttt{RIFIUTO: \{'A08', 'A25', 'A94', 'B19', 'B95', 'C10'\}} \\
            }
        }
    }
     \\ \\

    \theutterance \stepcounter{utterance}  
    & & \multicolumn{4}{p{0.6\linewidth}}{
        \cellcolor[rgb]{0.9,0.9,0.9}{
            \makecell[{{p{\linewidth}}}]{
                \texttt{\tiny{[P1$\langle$GM]}}
                \texttt{ARGOMENTAZIONE: \{"La tua proposta iniziale include oggetti con alta importanza, ma potrebbe essere utile considerare la distribuzione degli sforzi e le possibili alternative che potrebbero massimizzare l'importanza totale entro il limite. Alcuni oggetti nella tua proposta hanno sforzi considerevoli. Sarebbe interessante esplorare se ci sono combinazioni che includono oggetti con importanza simile ma con sforzi più contenuti, o se ci sono altri oggetti non inclusi nella tua proposta che potrebbero essere più vantaggiosi per entrambi."\}} \\
\\ 
\texttt{RIFIUTO: \{'A08', 'A25', 'A94', 'B19', 'B95', 'C10'\}} \\
            }
        }
    }
    & \\ \\

    \theutterance \stepcounter{utterance}  
    & \multicolumn{4}{p{0.6\linewidth}}{
        \cellcolor[rgb]{0.9,0.9,0.9}{
            \makecell[{{p{\linewidth}}}]{
                \texttt{\tiny{[P1$\rangle$GM]}}
                \texttt{RAGIONAMENTO STRATEGICO: \{'Il rifiuto della mia proposta iniziale suggerisce che l'altro partecipante sta cercando una combinazione più efficiente in termini di sforzo e importanza. Considererò quindi gli oggetti con sforzi più contenuti ma con un'importanza significativa. È anche possibile che l'altro partecipante abbia valutazioni diverse per gli oggetti e quindi potrei dovermi concentrare su oggetti che potrebbero essere meno controversi o che offrono un buon compromesso. La prossima proposta dovrà tenere conto di queste considerazioni per aumentare le possibilità di accettazione.'\}} \\
\\ 
\texttt{ARGOMENTAZIONE: \{"Capisco che la mia prima proposta potrebbe non essere stata ottimale in termini di equilibrio tra importanza e sforzo. Sono disposto a considerare alternative che potrebbero offrire un miglior compromesso. Alcuni oggetti che potrebbero essere presi in considerazione sono quelli con sforzi più bassi ma con importanza significativa per entrambe le parti. Sono aperto a discutere e trovare una soluzione che massimizzi l'importanza totale entro il limite di sforzo condiviso."\}} \\
\\ 
\texttt{PROPOSTA: \{'A87', 'B64', 'C79', 'B16', 'C53', 'A71'\}} \\
            }
        }
    }
    & & \\ \\

    \theutterance \stepcounter{utterance}  
    & & \multicolumn{4}{p{0.6\linewidth}}{
        \cellcolor[rgb]{0.9,0.9,0.9}{
            \makecell[{{p{\linewidth}}}]{
                \texttt{\tiny{[GM$\rangle$P2]}}
                \texttt{ARGOMENTAZIONE: \{"Capisco che la mia prima proposta potrebbe non essere stata ottimale in termini di equilibrio tra importanza e sforzo. Sono disposto a considerare alternative che potrebbero offrire un miglior compromesso. Alcuni oggetti che potrebbero essere presi in considerazione sono quelli con sforzi più bassi ma con importanza significativa per entrambe le parti. Sono aperto a discutere e trovare una soluzione che massimizzi l'importanza totale entro il limite di sforzo condiviso."\}} \\
\\ 
\texttt{PROPOSTA: \{'A87', 'B64', 'C79', 'B16', 'C53', 'A71'\}} \\
            }
        }
    }
    & \\ \\

    \theutterance \stepcounter{utterance}  
    & & & \multicolumn{4}{p{0.6\linewidth}}{
        \cellcolor[rgb]{0.9,0.9,0.9}{
            \makecell[{{p{\linewidth}}}]{
                \texttt{\tiny{[GM$\langle$P2]}}
                \texttt{RAGIONAMENTO STRATEGICO: \{"La risposta dell'altro giocatore mostra una volontà di negoziare e trovare un compromesso. La nuova proposta include oggetti con sforzi più bassi e potrebbe essere più equilibrata in termini di importanza e sforzo. Devo valutare l'importanza totale di questa proposta e considerare se gli oggetti selezionati sono quelli che massimizzano l'importanza per me. Inoltre, potrei considerare di proporre alternative che includono oggetti con alta importanza per me e che potrebbero essere più vantaggiosi per entrambi. È anche importante considerare come l'altro giocatore reagirà alle mie contro{-}proposte e come posso influenzare la sua percezione dell'importanza degli oggetti."\}} \\
\\ 
\texttt{ARGOMENTAZIONE: \{"La tua seconda proposta sembra essere più equilibrata in termini di sforzo, ma potremmo ancora esplorare altre opzioni per massimizzare l'importanza totale. Alcuni oggetti nella tua proposta, come 'A87' e 'B64', hanno un'importanza significativa, ma potremmo anche considerare oggetti come 'A39', 'A65' e 'C83' che potrebbero offrire un buon compromesso tra importanza e sforzo. Sarebbe utile discutere ulteriormente per capire quali oggetti sono veramente fondamentali per entrambi e trovare un accordo che tenga conto delle esigenze di entrambe le parti."\}} \\
\\ 
\texttt{PROPOSTA: \{'A39', 'A65', 'C83', 'B16', 'C53', 'A71', 'A28'\}} \\
            }
        }
    }
     \\ \\

    \theutterance \stepcounter{utterance}  
    & & \multicolumn{4}{p{0.6\linewidth}}{
        \cellcolor[rgb]{0.9,0.9,0.9}{
            \makecell[{{p{\linewidth}}}]{
                \texttt{\tiny{[P1$\langle$GM]}}
                \texttt{ARGOMENTAZIONE: \{"La tua seconda proposta sembra essere più equilibrata in termini di sforzo, ma potremmo ancora esplorare altre opzioni per massimizzare l'importanza totale. Alcuni oggetti nella tua proposta, come 'A87' e 'B64', hanno un'importanza significativa, ma potremmo anche considerare oggetti come 'A39', 'A65' e 'C83' che potrebbero offrire un buon compromesso tra importanza e sforzo. Sarebbe utile discutere ulteriormente per capire quali oggetti sono veramente fondamentali per entrambi e trovare un accordo che tenga conto delle esigenze di entrambe le parti."\}} \\
\\ 
\texttt{PROPOSTA: \{'A39', 'A65', 'C83', 'B16', 'C53', 'A71', 'A28'\}} \\
            }
        }
    }
    & \\ \\

    \theutterance \stepcounter{utterance}  
    & \multicolumn{4}{p{0.6\linewidth}}{
        \cellcolor[rgb]{0.9,0.9,0.9}{
            \makecell[{{p{\linewidth}}}]{
                \texttt{\tiny{[P1$\rangle$GM]}}
                \texttt{RAGIONAMENTO STRATEGICO: \{'La proposta dell'altro partecipante introduce nuovi oggetti che potrebbero offrire un buon equilibrio tra importanza e sforzo. È interessante notare l'inclusione di 'A39', 'A65' e 'C83', che potrebbero essere più efficienti in termini di sforzo rispetto ad alcuni oggetti della mia precedente proposta. Tuttavia, devo valutare attentamente l'importanza di questi oggetti rispetto ai miei obiettivi e considerare se la loro inclusione massimizzi effettivamente l'importanza totale. Potrei dover negoziare ulteriormente per includere oggetti che sono fondamentali per me, magari proponendo scambi o sostituzioni che beneficino entrambe le parti.'\}} \\
\\ 
\texttt{ARGOMENTAZIONE: \{"La tua proposta presenta alcuni oggetti interessanti che potrebbero essere più efficienti in termini di sforzo. Tuttavia, vorrei discutere ulteriormente sull'inclusione di 'A87' e 'B64' nella mia precedente proposta, poiché questi oggetti hanno un'importanza significativa per me. Sarebbe utile esplorare la possibilità di sostituire alcuni oggetti nella tua proposta con quelli che sono più cruciali per me, trovando un equilibrio che massimizzi l'importanza totale per entrambi."\}} \\
\\ 
\texttt{RIFIUTO: \{'A39', 'A65', 'C83', 'B16', 'C53', 'A71', 'A28'\}} \\
            }
        }
    }
    & & \\ \\

    \theutterance \stepcounter{utterance}  
    & & \multicolumn{4}{p{0.6\linewidth}}{
        \cellcolor[rgb]{0.9,0.9,0.9}{
            \makecell[{{p{\linewidth}}}]{
                \texttt{\tiny{[GM$\rangle$P2]}}
                \texttt{ARGOMENTAZIONE: \{"La tua proposta presenta alcuni oggetti interessanti che potrebbero essere più efficienti in termini di sforzo. Tuttavia, vorrei discutere ulteriormente sull'inclusione di 'A87' e 'B64' nella mia precedente proposta, poiché questi oggetti hanno un'importanza significativa per me. Sarebbe utile esplorare la possibilità di sostituire alcuni oggetti nella tua proposta con quelli che sono più cruciali per me, trovando un equilibrio che massimizzi l'importanza totale per entrambi."\}} \\
\\ 
\texttt{RIFIUTO: \{'A39', 'A65', 'C83', 'B16', 'C53', 'A71', 'A28'\}} \\
            }
        }
    }
    & \\ \\

    \theutterance \stepcounter{utterance}  
    & & & \multicolumn{4}{p{0.6\linewidth}}{
        \cellcolor[rgb]{0.9,0.9,0.9}{
            \makecell[{{p{\linewidth}}}]{
                \texttt{\tiny{[GM$\langle$P2]}}
                \texttt{RAGIONAMENTO STRATEGICO: \{"L'altro giocatore ha rifiutato la mia proposta e ha espresso interesse per la sostituzione di alcuni oggetti con 'A87' e 'B64', che sono importanti per lui. Questo indica che questi oggetti hanno un valore elevato per l'altro giocatore e che potrebbero essere utili come oggetti di scambio. Devo considerare se accettare di includere 'A87' e 'B64' nella proposta in cambio di altri oggetti che sono più importanti per me. Inoltre, potrei cercare di capire quali altri oggetti sono importanti per l'altro giocatore e utilizzare questa informazione per negoziare un accordo più favorevole."\}} \\
\\ 
\texttt{ARGOMENTAZIONE: \{"Capisco che 'A87' e 'B64' sono importanti per te, e sono disposto a considerare la loro inclusione in una proposta. Tuttavia, vorrei anche assicurarmi che gli oggetti più importanti per me siano inclusi. Una possibile soluzione potrebbe essere trovare un equilibrio tra le nostre esigenze, magari sostituendo alcuni oggetti con sforzo più alto con altri che hanno un impatto significativo sulla nostra importanza totale. Sarebbe utile discutere ulteriormente per capire quali oggetti sono veramente fondamentali per entrambi e trovare un accordo che tenga conto delle esigenze di entrambe le parti."\}} \\
\\ 
\texttt{PROPOSTA: \{'A87', 'B64', 'C83', 'B16', 'C53', 'A71', 'A39'\}} \\
            }
        }
    }
     \\ \\

    \theutterance \stepcounter{utterance}  
    & & \multicolumn{4}{p{0.6\linewidth}}{
        \cellcolor[rgb]{0.9,0.9,0.9}{
            \makecell[{{p{\linewidth}}}]{
                \texttt{\tiny{[P1$\langle$GM]}}
                \texttt{ARGOMENTAZIONE: \{"Capisco che 'A87' e 'B64' sono importanti per te, e sono disposto a considerare la loro inclusione in una proposta. Tuttavia, vorrei anche assicurarmi che gli oggetti più importanti per me siano inclusi. Una possibile soluzione potrebbe essere trovare un equilibrio tra le nostre esigenze, magari sostituendo alcuni oggetti con sforzo più alto con altri che hanno un impatto significativo sulla nostra importanza totale. Sarebbe utile discutere ulteriormente per capire quali oggetti sono veramente fondamentali per entrambi e trovare un accordo che tenga conto delle esigenze di entrambe le parti."\}} \\
\\ 
\texttt{PROPOSTA: \{'A87', 'B64', 'C83', 'B16', 'C53', 'A71', 'A39'\}} \\
            }
        }
    }
    & \\ \\

    \theutterance \stepcounter{utterance}  
    & \multicolumn{4}{p{0.6\linewidth}}{
        \cellcolor[rgb]{0.9,0.9,0.9}{
            \makecell[{{p{\linewidth}}}]{
                \texttt{\tiny{[P1$\rangle$GM]}}
                \texttt{RAGIONAMENTO STRATEGICO: \{'La nuova proposta dell'altro partecipante sembra essere più incline a includere oggetti importanti per me, come 'A87' e 'B64'. L'inclusione di 'C83', 'A39' e altri oggetti con sforzo relativamente basso potrebbe essere un buon compromesso. Tuttavia, devo valutare attentamente se questa combinazione massimizza effettivamente l'importanza totale per me. Considererò attentamente l'importanza di ogni oggetto in questa proposta e valuterò se ci sono ulteriori opportunità per migliorare l'accordo. Potrei accettare questa proposta se ritengo che rappresenti il miglior compromesso possibile entro il limite di sforzo.'\}} \\
\\ 
\texttt{ARGOMENTAZIONE: \{"La tua ultima proposta sembra essere un passo avanti significativo verso un accordo. L'inclusione di 'A87' e 'B64' è cruciale per me, e gli altri oggetti scelti sembrano offrire un buon equilibrio tra importanza e sforzo. Sono disposto a considerare ulteriori ajustamenti minori se necessario, ma in generale, questa proposta sembra essere una base solida per un accordo. Sono pronto ad accettare questa proposta se ritengo che non ci siano alternative significativamente migliori."\}} \\
\\ 
\texttt{ACCORDO: \{'A87', 'B64', 'C83', 'B16', 'C53', 'A71', 'A39'\}} \\
            }
        }
    }
    & & \\ \\

    \theutterance \stepcounter{utterance}  
    & & & \multicolumn{2}{p{0.3\linewidth}}{
        \cellcolor[rgb]{0.9,0.9,0.9}{
            \makecell[{{p{\linewidth}}}]{
                \texttt{\tiny{[GM$|$GM]}}
                \texttt{game successful} \\
            }
        }
    }
    & & \\ \\

    \theutterance \stepcounter{utterance}  
    & & & \multicolumn{2}{p{0.3\linewidth}}{
        \cellcolor[rgb]{0.9,0.9,0.9}{
            \makecell[{{p{\linewidth}}}]{
                \texttt{\tiny{[GM$|$GM]}}
                \texttt{end game} \\
            }
        }
    }
    & & \\ \\

\end{supertabular}
}

\end{document}
